\chapter{Background and Literature Review}
\label{chap:background}

% TODO: bridging paragraph introducing the three pillars surveyed in this
%       chapter: (i) soft-robot modelling, (ii) learning-based IK, (iii)
%       simulation and sim-to-real transfer.

\section{Pneumatic soft manipulators}
\label{sec:bg_pneumatic}

The two datasets that underpin this thesis were both recorded on pneumatic soft manipulators. This section describes the design principles of that class of robot (Section~\ref{sec:bg_pneumatic_design}), the I-Support platform from which Dataset 1 was recorded (Section~\ref{sec:bg_pneumatic_isupport}), and the smaller family of three-chamber and three-valve platforms that is relevant to Dataset 2 and to the Ma et al.\ \cite{ma2022bp} benchmark (Section~\ref{sec:bg_pneumatic_others}).

\subsection{Design and actuation principles}
\label{sec:bg_pneumatic_design}

A pneumatic soft manipulator is a continuum robot whose structure deforms under fluid pressure rather than by rotating rigid links at discrete joints. The canonical design is a silicone cylinder with several sealed chambers distributed around its neutral axis; pressurising a subset of the chambers lengthens them asymmetrically relative to the others and bends the cylinder away from the active chambers. When all chambers are pressurised equally the cylinder elongates uniformly, and when only one is pressurised the cylinder bends in the plane defined by that chamber and the axis. The most common arrangement uses three chambers placed at $120^{\circ}$ around the axis, which allows bending in any direction in the plane perpendicular to the axis using only three independent pressures. Commercial and research platforms add fibre reinforcement around the chambers---typically a braided polyester or polyethylene terephthalate sheath---to constrain radial expansion and to channel pressure into axial bending or elongation.

The appeal of pneumatic actuation for soft manipulators is the combination of intrinsic mechanical compliance, safety for unstructured contact with humans, and simplicity at the drive end: pressurised air is a cheap and widely available power source, valves are small, and the actuator itself contains no internal friction surfaces. The cost is that the relationship between the applied pressure and the resulting deformation is strongly nonlinear and history-dependent. The silicone matrix exhibits viscoelastic creep on the order of seconds, so the deformation reached under a given pressure depends on the history of that pressure and not only on its current value. Manufacturing variability is also substantial: lateral displacements of $25$--$30$\,mm have been reported for otherwise identical I-Support chambers as a result of manual assembly tolerances \cite{manti2016soft}. These properties together mean that rigid-robot inverse-kinematics methods do not transfer directly to pneumatic soft arms, and they are the main reason the soft-robotics community has pursued learning-based control strategies \cite{kim2021mlsoft}.

\subsection{The I-Support platform}
\label{sec:bg_pneumatic_isupport}

The I-Support manipulator was developed by Manti et al.\ \cite{manti2016soft} at the BioRobotics Institute of Scuola Superiore Sant'Anna, originally as an assistive robot for the showering of elderly users. A single I-Support module is a silicone cylinder of $205$\,mm rest length and $60$\,mm diameter, with three McKibben-type fluidic chambers arranged at $120^{\circ}$ around the central axis and offset approximately $12$\,mm from it. The module is reinforced with a braided polyethylene terephthalate sheath that is thermally set to a bellows shape, which concentrates pressure-induced deformation into axial bending and elongation rather than radial ballooning. Each chamber can be safely pressurised up to $1.2$\,bar, at which a single-chamber actuation bends the module tip by approximately $118^{\circ}$ from its rest pose and a symmetric three-chamber actuation elongates the module by approximately $85\%$ of its rest length. These two behavioural anchors are the reference values used in the PyElastica calibration of Section~\ref{sec:sim_setup}.

Dataset 1 of this thesis was recorded on the two-module extension of the I-Support platform introduced by Bianchi et al.\ \cite{bianchi2023learning}. The two modules are stacked in series, rotated by $60^{\circ}$ relative to each other so that the six chambers together cover the full actuation circle, and mounted vertically with the end-effector pointing downwards so that gravity acts along the straight-pose axis rather than perpendicular to it. The assembled arm has a rest length of approximately $410$\,mm. These design choices together make the two-module I-Support a convenient reference platform for data-driven inverse-kinematics work on pneumatic soft robots: the actuation is pressure-only with no tendons or cables, the two-module geometry provides the six-chamber over-parameterisation that motivates multi-valued inverse methods (Section~\ref{sec:disc_interpretation}), and marker-based optical tracking of the base, mid-section, and tip provides a spatial sample of the internal configuration along the arm.

\subsection{Other pneumatic platforms relevant to this work}
\label{sec:bg_pneumatic_others}

Two further pneumatic platforms are relevant to this thesis. The first is the three-chamber single-segment arm of Ma et al.\ \cite{ma2022bp}, a purpose-built test rig developed in 2022 to evaluate a back-propagation neural network for inverse-kinematics prediction. Ma et al.'s arm uses three silicone chambers at $120^{\circ}$ around the axis without a second stage, with documented actuation units in bar and a documented straight-pose origin. Its single-segment three-chamber topology matches the geometric assumptions of the piecewise-constant-curvature model that Section~\ref{sec:bg_ik_analytical} reviews below, and the sub-$5\%$ pressure-space MAPE that Ma et al.\ reported on a $216$-sample dataset is a natural benchmark for data-driven pneumatic inverse kinematics; it is used as a scale marker for the Dataset 2 comparisons of Chapter~\ref{chap:results}.

The second is the undocumented three-valve soft arm from which Dataset 2 was recorded. The dataset was supplied by the project supervisor and is not otherwise in the public literature. From the dataset alone it is known that actuation is driven by three valves on an integer-valued command range $[0, 2500]$, that tip positions were recorded in inches and converted to millimetres, and that the observed workspace is a shallow disc of approximately $100$\,mm lateral extent and $65$\,mm axial extent (Section~\ref{sec:ds2}). The absence of documented geometry, material, and actuation calibration places Dataset 2 at the opposite extreme from Ma et al.'s platform along the documentation axis, and several of the experimental limitations discussed in Chapter~\ref{chap:discussion} trace back directly to this asymmetry.

Beyond these two specific platforms, the last decade has produced a wide variety of pneumatic soft-manipulator designs, including single-chamber bending actuators, modular multi-chamber arms, inflatable robotic trunks, and fibre-reinforced continuum manipulators; Kim et al.\ \cite{kim2021mlsoft} and Alessi et al.\ \cite{alessi2024rodreview} provide recent surveys. The present thesis draws on only two datasets, chosen because they expose complementary design properties---a two-module over-actuated pneumatic recording in Dataset 1 and an under-documented three-valve recording in Dataset 2---that together cover the actuation-redundancy regimes that the method-selection discussion of Chapter~\ref{chap:discussion} depends on.

\section{Kinematic modelling of soft robots}
\label{sec:bg_kinematics}

The inverse-kinematics methods compared in this thesis rely on an underlying model of how actuation maps to arm shape and tip position. This section reviews the three classes of kinematic model relevant to continuum soft robots: the piecewise-constant-curvature (PCC) simplification that dominates the analytical literature (Section~\ref{sec:bg_kin_pcc}), its multi-segment composition for arms with more than one module (Section~\ref{sec:bg_kin_multi}), and the more general Cosserat-rod and finite-element formulations that trade analytical tractability for physical fidelity (Section~\ref{sec:bg_kin_cosserat}).

\subsection{Piecewise constant curvature}
\label{sec:bg_kin_pcc}

The piecewise-constant-curvature (PCC) model approximates a soft manipulator's backbone as a sequence of circular arcs, each parameterised by an arc length $L$, a bending angle $\theta$ (equivalently a curvature $\kappa = \theta / L$), and an azimuth $\varphi$ that selects the plane of the bend. A single segment therefore has three configuration-space degrees of freedom, and the forward map from configuration to tip position takes a standard closed form that Webster and Jones \cite{webster2010review} derived in the most widely cited survey of continuum-robot kinematics. The appeal of PCC is its tractability: the forward kinematics is a closed-form expression in $(L, \theta, \varphi)$, the Jacobian can be written analytically, and for a single segment the inverse kinematics reduces to a system of three equations in three unknowns that can be solved in microseconds. These properties have made PCC the standard analytical baseline for pneumatic and cable-driven continuum manipulators through most of the last decade.

The tractability comes at the cost of a restrictive modelling assumption. PCC requires that the curvature is constant along every point of a segment, which in turn requires that the forces and moments acting along the segment are uniform. This assumption is approximately met on short, stiff, tendon-driven arms in which gravity and external contact are small compared with the actuation-induced bending moment. It breaks down on longer pneumatic arms subject to gravity, on arms with distributed mass or stiffness variation along the backbone, and on arms where the fluid-induced bending moment itself is not uniform. Published reviews routinely identify PCC's constant-curvature assumption as the principal source of residual error when the predicted and measured tip positions are compared on real hardware \cite{webster2010review, alessi2024rodreview}. The Dataset 2 results of Chapter~\ref{chap:results} provide one further empirical instance: the single-segment PCC baseline reaches a $22$\,mm forward residual on a workspace of only $100$\,mm because an elongation coupling that a constant-$L$ PCC cannot represent dominates the error (Section~\ref{sec:res_pcc_single}).

\subsection{Multi-segment extensions}
\label{sec:bg_kin_multi}

Multi-segment soft manipulators such as the two-module I-Support of Section~\ref{sec:bg_pneumatic_isupport} are modelled by chaining single-segment PCC transforms in series. Sokolov et al.\ \cite{sokolov2024twoseg} present an explicit two-segment derivation with a $60^{\circ}$ inter-segment rotation that corresponds to the I-Support geometry exactly, and this is the formulation adopted for the two-segment PCC baseline of Section~\ref{sec:pcc_multi}. The forward map remains closed-form under composition, but the inverse map loses its single-valued character: a two-segment arm typically has six or more actuation degrees of freedom that map into a three-dimensional tip position, and the resulting inverse is a relation rather than a function.

Two main strategies have emerged for solving the multi-segment inverse. The first is an optimisation-based search over the configuration space, constrained to physically plausible per-segment bounds; Keyvanara et al.\ \cite{keyvanara2023geometric} use sequential least-squares quadratic programming for this purpose, with a trajectory-continuity warm-start from the previous target that this thesis adopts in Section~\ref{sec:pcc_multi}. The second is a Jacobian-based iteration on the analytically tractable forward map, as developed for cable-driven continuum arms by Giorelli et al.\ \cite{giorelli2015nn}. Both approaches recover the configuration of the arm to within numerical tolerance, but they leave open how that configuration maps to the measurable actuation command; that step is straightforward when the actuation units are documented and calibrated, and it collapses under ordinary least squares when many distinct configurations share nearly the same chamber-length vector. The latter failure mode is the one observed on Dataset 1 in Section~\ref{sec:res_pcc_multi}.

\subsection{Beyond PCC: Cosserat rod and finite-element approaches}
\label{sec:bg_kin_cosserat}

When the constant-curvature assumption of PCC is violated, Cosserat-rod theory provides a more general framework that models the backbone as a continuous elastic body with cross-section orientation and position evolving along an arc-length parameter. The Cosserat rod captures nonconstant curvature under distributed loads, couples bending with elongation and torsion through the material's elastic moduli, and produces configurations that PCC cannot reach by construction. Alessi et al.\ \cite{alessi2024rodreview} review the use of rod models in continuum and soft-robot control and distinguish four main rod theories that differ in the deformation modes they represent and in how they are discretised for numerical integration. The PyElastica library of Gazzola et al.\ \cite{gazzola2018pyelastica}, which this thesis uses for the simulator of Section~\ref{sec:sim_setup}, implements the special-Cosserat-rod formulation with a position-Verlet integrator that scales linearly in the number of discretisation elements.

Higher-fidelity modelling is possible with fully three-dimensional finite-element approaches, in which the continuum is discretised into volumetric elements rather than a one-dimensional rod. Finite-element methods offer the richest physical accuracy---including contact, material nonlinearity, and stress distribution in cross-section---at substantially higher computational cost and with far more complex parameter calibration than either PCC or Cosserat-rod models. Both Cosserat-rod and finite-element models share the property that their inverse is expensive to compute directly, which is why learning-based methods that fit a fast neural surrogate to expensive simulator trajectories---the sim-to-real approach of Fang et al.\ \cite{fang2022jacobian} and of the PyElastica pipeline evaluated in Chapter~\ref{chap:results}---have become a standard route around the modelling-versus-tractability trade-off.

The practical outcome of the trade-off is that PCC remains the standard analytical baseline in soft-robot inverse-kinematics work despite its limitations: it is fast, interpretable, and accurate enough for short single-segment arms where its assumptions approximately hold. Cosserat-rod models enter when gravity, elongation, or nonconstant loading cannot be neglected, and finite-element methods are reserved for design-space studies in which simulation time is not a control-loop constraint. This thesis uses PCC as a classical reference in Section~\ref{sec:pcc_single}, the Sokolov two-segment composition in Section~\ref{sec:pcc_multi}, and a PyElastica Cosserat rod as the basis of the sim-to-real study in Section~\ref{sec:sim_setup}, covering the three regimes of the trade-off in a single experimental programme.

\section{Inverse Kinematics Methods}

\subsection{Analytical IK}
% TODO: geometric / trigonometric inversion of the PCC model; limitations
%       when the real robot deviates from the constant-curvature assumption.

\subsection{Feed-Forward Neural IK}
% TODO: MLP regressors from target pose to actuation; Giorelli et al. and
%       Ma et al. (2022) as canonical references.

\subsection{Recurrent and Temporal IK}
% TODO: LSTM / GRU for hysteresis and history-dependent deformation.

\subsection{Jacobian-Learned IK}
% TODO: Fang et al. (2022) damped least-squares scheme with a learned
%       Jacobian; pros vs. direct regression.

\section{Simulation and Sim-to-Real Transfer}

\subsection{PyElastica for Soft-Robot Simulation}
% TODO: Cosserat-rod formulation, pneumatic actuation modelling.

\subsection{Sim-to-Real Strategies}
% TODO: domain randomisation, residual learning, fine-tuning; applicability
%       to pneumatic soft arms.

\section{Evaluation Metrics and Benchmarks}

\subsection{Pose Error Metrics}
% TODO: MAE / RMSE / MAPE of tip position; percentage of workspace diameter.

\subsection{Trajectory Tracking Metrics}
% TODO: path-following error, smoothness, actuation effort.

\subsection{Published Baselines}
% TODO: headline numbers from Ma et al. (2022), Giorelli et al., Fang et al.
%       as the bar to clear.

\section{Summary and Research Gap}
% TODO: state the gap this thesis fills: a head-to-head comparison of
%       analytical, MLP, LSTM and Jacobian-learned IK on identical data,
%       plus a PyElastica sim-to-real study.
